\chapter{The simulation engine}\label{ch:engine}

The engine (\file{linkstudio/core}) is about one thousand lines of Python and contains no
communication theory. It provides five abstractions:
\begin{itemize}
  \item \emph{signals} that flow along edges;
  \item \emph{blocks} that transform them;
  \item \emph{graphs} that connect blocks;
  \item \emph{simulations} that run graphs batch after batch;
  \item \emph{sweeps} that repeat simulations over a parameter grid.
\end{itemize}
This chapter describes each of them and the rules they enforce.

% ------------------------------------------------------------------
\section{Signals}\label{sec:signal}
A \code{Signal} (\file{core/signal.py}) is an immutable record with four fields:
\begin{description}[font=\ttfamily\color{accent}]
  \item[data] a NumPy array. Axis 0 is always the \emph{batch} (independent frames simulated in
  parallel) and the last axis is time, bit index or subcarrier index. Multi-dimensional kinds
  add axes in between, as in \cref{fig:signal-shapes}.
  \item[kind] the semantic type of the data, a flag of the \code{Kind} enumeration
  (\cref{tab:kinds}).
  \item[fs] the sample or symbol rate in Hz, when it is meaningful.
  \item[meta] a dictionary of metadata that travels downstream with the data (\cref{sec:metadata}).
\end{description}

\begin{figure}[htbp]
\centering
\begin{tikzpicture}[x=1cm,y=1cm]
  % 2D
  \begin{scope}
    \foreach \r in {0,...,3} \foreach \c in {0,...,7}
      \draw[fill=ksym!10,draw=ksym!60] (\c*0.4,-\r*0.4) rectangle ++(0.4,0.4);
    \draw[->] (-0.2,0.6) -- node[above,lbl]{time (last axis)} (3.4,0.6);
    \draw[->] (-0.25,0.4) -- node[left,lbl,align=right]{batch\\(axis 0)} (-0.25,-1.2);
    \node[lbl,below] at (1.6,-1.35) {\code{(B, n)}: bits, symbols, samples};
  \end{scope}
  % 3D streams
  \begin{scope}[xshift=4.9cm]
    \foreach \p in {1,0} {
      \foreach \r in {0,...,3} \foreach \c in {0,...,6}
        \draw[fill=ksamp!10,draw=ksamp!60] (\c*0.4+\p*0.25,-\r*0.4+\p*0.25) rectangle ++(0.4,0.4);}
    \node[lbl] at (3.45,0.85) {stream / polarization};
    \node[lbl,below] at (1.6,-1.35) {\code{(B, P, n)}: MIMO, dual pol.};
  \end{scope}
  % grid
  \begin{scope}[xshift=9.6cm]
    \foreach \p in {2,1,0} {
      \foreach \r in {0,...,3} \foreach \c in {0,...,5}
        \draw[fill=kgrid!12,draw=kgrid!80] (\c*0.4+\p*0.22,-\r*0.4+\p*0.22) rectangle ++(0.4,0.4);}
    \node[lbl,align=center] at (1.5,1.15) {OFDM symbols $\times$ subcarriers};
    \node[lbl,below] at (1.4,-1.35) {\code{(B, S, N_FFT)}: resource grid};
  \end{scope}
\end{tikzpicture}
\caption{Array conventions. Every block may assume that axis 0 is the batch and the last axis is
time (or frequency for grids).}
\label{fig:signal-shapes}
\end{figure}

\begin{table}[htbp]
\centering\small
\caption{Signal kinds. Kinds are bit flags, so a port can accept several of them
(\code{Kind.BITS | Kind.LLR}). The colours are those used by the GUI.}
\label{tab:kinds}
\begin{tabular}{@{}lllp{6.8cm}@{}}
\toprule
Kind & & dtype & Meaning\\
\midrule
\kind{BITS}    & \textcolor{kbits}{\rule{8pt}{8pt}} & \code{uint8} & hard bits in $\{0,1\}$\\
\kind{LLR}     & \textcolor{kllr}{\rule{8pt}{8pt}}  & \code{float64} & log-likelihood ratios, $L=\ln P(b{=}0)/P(b{=}1)$ (positive means 0)\\
\kind{SYMBOLS} & \textcolor{ksym}{\rule{8pt}{8pt}}  & \code{complex128} & one sample per symbol (constellation domain)\\
\kind{SAMPLES} & \textcolor{ksamp}{\rule{8pt}{8pt}} & \code{complex128}/\code{float64} & oversampled waveform; optical fields in $\sqrt{\mathrm{W}}$\\
\kind{GRID}    & \textcolor{kgrid}{\rule{8pt}{8pt}} & \code{complex128} & OFDM resource grid (frequency domain)\\
\kind{CSI}     & \textcolor{kcsi}{\rule{8pt}{8pt}}  & \code{complex128} & channel state information; the form is given by \meta{csi_type}\\
\kind{FLAGS}   & \textcolor{black!50}{\rule{8pt}{8pt}} & \code{bool} & one flag per frame (e.g. CRC check)\\
\bottomrule
\end{tabular}
\end{table}

Signals are never modified in place. A block derives a new signal with
\code{x.evolve(new_data, key=value)}. This keeps the input's kind, rate and metadata and
merges the new keys, so information added upstream is preserved unless a block deliberately
overrides it.

% ------------------------------------------------------------------
\section{Blocks}\label{sec:blocks}
A block (\file{core/block.py}) declares three things:
\begin{itemize}
  \item a \code{Params} model, a Pydantic class whose fields carry types, defaults, bounds
  and descriptions;
  \item named, typed \code{inputs} and \code{outputs} (\code{Port} objects);
  \item a \code{process(inputs, ctx)} method mapping one batch of input signals to output
  signals.
\end{itemize}
Measurement blocks also override \code{metrics()}.

Ports carry the accepted (or produced) kinds, a \code{required} flag for inputs and, for
outputs, an optional \code{same_as} reference: \code{Port.like("in")} produces the same kind
as input \code{in}. The GUI uses this information to colour handles and to refuse invalid
connections, and the graph validator uses it to propagate kinds (\cref{sec:validation}).

\subsection{Life cycle}
\begin{figure}[htbp]
\centering
\begin{tikzpicture}[node distance=6mm and 8mm]
  \node[blockbox] (init) {\code{__init__}\\validate Params};
  \node[blockbox,right=of init] (setup) {\code{setup()}\\derived state};
  \node[blockbox,right=of setup] (reset) {\code{reset(seed, path)}\\RNG + \code{on_reset()}};
  \node[greenbox,right=of reset] (proc) {\code{process()}\\one batch};
  \node[blockbox,below=9mm of proc] (met) {\code{metrics()}};
  \node[warmbox,below=9mm of setup] (setp) {\code{set_param()}};
  \draw[sig] (init) -- (setup);
  \draw[sig] (setup) -- (reset);
  \draw[sig] (reset) -- (proc);
  \draw[sig] (proc.north east) to[out=60,in=0,looseness=4] node[right,lbl]{next batch} (proc.east);
  \draw[sig] (proc) -- node[right,lbl]{any time} (met);
  \draw[sig] (setp) -- node[left,lbl]{rebuild} (setup);
\end{tikzpicture}
\caption{Block life cycle. \code{setup()} rebuilds derived state (filter taps, code matrices,
trellises) after construction and after every parameter change. \code{on_reset()} clears
accumulators before a run.}
\label{fig:lifecycle}
\end{figure}

\subsection{Reproducible randomness}
Every block owns a private NumPy generator. On reset it is re-seeded from the master seed of
the run \emph{and} the block's hierarchical path:
\codeexcerpt{block_reset}{Seeding the random generator of a block.}
Because the CRC-32 of the path enters the seed sequence, two \blk{awgn} blocks named
differently draw independent noise. The same graph with the same seed produces bit-identical
results, whatever the execution order of unrelated blocks and whether the run starts from the
GUI, the CLI or a script.

% ------------------------------------------------------------------
\section{The registry}
\code{@register_block("type_name", category=...)} adds a class to a global registry
(\file{core/registry.py}). Registered blocks can be created by name (\code{create_block},
\code{Graph.create}) and are described by \code{block_catalogue()}. The catalogue contains
the documentation, the ports and the JSON Schema of the parameters. The GUI builds its palette
and its parameter forms from this catalogue, and \cref{app:blocks} of this manual is generated
from it.

% ------------------------------------------------------------------
\section{Graphs}\label{sec:graphs}
A \code{Graph} holds named blocks and edges. Edges are stored as a map
$\text{destination input}\mapsto\text{source output}$, so every input has at most one driver,
while an output may fan out to any number of inputs. Ports are addressed as
\code{"block.port"}; omitting the port selects the block's first declared port, which makes
\code{g.chain("src", "map", "awgn")} a compact way to wire a pipeline. Parameters are
addressed as \code{"block.param"}, and \code{"/"} descends into subsystems
(\code{"rx/demap.method"}).

\subsection{Validation}\label{sec:validation}
\code{Graph.validate()} enforces three rules and returns the resolved kind of every output:
\begin{enumerate}
  \item every \emph{required} input is driven, either by an edge or, inside a subsystem, by an
  external input;
  \item the graph is acyclic;
  \item along every edge the produced kind intersects the accepted kinds.
\end{enumerate}
Kinds are resolved in topological order. An output declared with \code{same_as} takes the
kind actually delivered to the referenced input, intersected with the kinds that input accepts.
An \blk{interleaver} fed with LLRs therefore produces \kind{LLR}, and the validator can reject
an LLR output connected to a block that only accepts symbols, before anything is executed.

The topological order is computed with Kahn's algorithm:
\codeexcerpt{graph_toposort}{Topological sort of the block graph.}
Ties are broken by insertion order, so the execution order is deterministic. Feedback loops
(PLLs, adaptive filters, decision feedback) live \emph{inside} blocks, where they can run at the
sample rate. The graph itself carries one batch at a time from sources to sinks.

\begin{figure}[htbp]
\centering
\begin{tikzpicture}[node distance=6mm and 6mm]
  \node[blockbox] (src) {src\\\scriptsize\color{accent2}1};
  \node[blockbox,right=of src] (enc) {enc\\\scriptsize\color{accent2}2};
  \node[blockbox,right=of enc] (map) {map\\\scriptsize\color{accent2}3};
  \node[warmbox,right=of map] (ch) {ch\\\scriptsize\color{accent2}4};
  \node[blockbox,right=of ch] (eq) {eq\\\scriptsize\color{accent2}5};
  \node[blockbox,right=of eq] (dec) {dec\\\scriptsize\color{accent2}6};
  \node[greenbox,right=of dec] (ber) {ber\\\scriptsize\color{accent2}7};
  \draw[sig] (src)--(enc); \draw[sig] (enc)--(map); \draw[sig] (map)--(ch);
  \draw[sig] (ch)--node[above,lbl]{out}(eq); \draw[sig] (eq)--(dec); \draw[sig] (dec)--node[above,lbl]{rx}(ber);
  \draw[sig,kcsi] (ch.south) -- ++(0,-4mm) -| node[pos=0.25,below,lbl]{csi} (eq.south);
  \draw[sig,kbits] (src.north) -- ++(0,5mm) -| node[pos=0.25,above,lbl]{tx} (ber.north);
\end{tikzpicture}
\caption{A graph with fan-out (\code{src.out} drives \code{enc.in} and \code{ber.tx}) and a
side channel (\code{ch.csi}). Red numbers are the execution order.}
\end{figure}

\subsection{Execution}
\code{Graph.execute(ctx)} walks the blocks in order, gathers their inputs from the signals
already produced, checks the delivered kinds at run time and calls \code{process}:
\codeexcerpt{graph_execute}{The core of the executor.}
All produced signals are returned keyed by \code{"block.port"}; internal signals of
subsystems appear as \code{"sub/block.port"}. These keys are what the GUI probes and what
\code{Simulation.run(capture=...)} returns. Exceptions raised inside a block are re-raised as
\code{BlockError} with the block name, so the user sees which block failed and why.

% ------------------------------------------------------------------
\section{Subsystems}\label{sec:subsystems}
A \code{Subsystem} (\file{core/composite.py}) wraps a whole graph as a single block.
Its ports map to chosen inner ports, and its output kinds are those resolved by validating the
inner graph. A transmitter or a receiver can thus be built once and reused in many scenarios;
the same mechanism will let a future network layer treat a whole PHY link as one node.
Parameters and metrics of inner blocks are reached with the \code{"/"} separator:
\code{rx/dec.max_iter}, \code{rx/dec/avg_iterations}.

\begin{figure}[htbp]
\centering
\begin{tikzpicture}[node distance=5mm and 7mm]
  \node[blockbox] (src) {src};
  \node[draw=accent4,dashed,rounded corners,inner sep=7pt,right=12mm of src,
        label={[lbl,accent4]above:subsystem \code{tx} (file \file{coded_qam_tx.yaml})}] (tx) {
    \tikz[node distance=5mm]{\node[blockbox,minimum width=10mm] (e) {enc}; \node[blockbox,minimum width=10mm,right=of e] (m) {map}; \draw[sig] (e)--(m);}};
  \node[warmbox,right=12mm of tx] (ch) {awgn};
  \node[draw=accent4,dashed,rounded corners,inner sep=7pt,right=12mm of ch,
        label={[lbl,accent4]above:subsystem \code{rx}}] (rx) {
    \tikz[node distance=5mm]{\node[blockbox,minimum width=10mm] (d) {demap}; \node[blockbox,minimum width=10mm,right=of d] (c) {dec}; \draw[sig] (d)--(c);}};
  \draw[sig] (src) -- node[above,lbl]{bits} (tx);
  \draw[sig] (tx) -- node[above,lbl]{symbols} (ch);
  \draw[sig] (ch) -- (rx);
\end{tikzpicture}
\caption{Subsystems nest graphs; from outside they behave as ordinary blocks.}
\end{figure}

% ------------------------------------------------------------------
\section{Simulations and stop criteria}\label{sec:simulation}
A \code{Simulation} (\file{core/simulation.py}) runs a graph batch after batch. One batch
contains \code{batch_size} independent frames, which NumPy processes as a single vectorized
operation, so a larger batch amortizes the Python overhead. Meter blocks accumulate counts
across batches. The run stops at the first satisfied condition of a \code{StopCriteria}:
\codeexcerpt{stop_check}{Stop criteria are checked after every batch.}
The classic rule for error-rate estimation is \code{StopCriteria.errors("ber.bit_errors", 100)}:
stop once 100 errors have been counted. \Cref{sec:mc} shows why this gives a relative
accuracy of about $\pm20\,\%$ independent of the BER itself. \code{run} also accepts an
\code{on_batch} callback, which the web server uses to report progress and to cancel jobs.

% ------------------------------------------------------------------
\section{Sweeps}\label{sec:sweeps}
A \code{Sweep} (\file{core/sweep.py}) runs a fresh simulation at every point of a parameter
grid. In \code{grid} mode the points are the Cartesian product of the value lists; in
\code{zip} mode the lists are paired element-wise. It returns a table with one record per
point, which can be filtered, printed or written to CSV.
\codeexcerpt{sweep_point}{Each sweep point sets the parameters and runs a new simulation.}

\begin{keypoint}[Common random numbers]
By default every point reuses the same seed (\code{same_seed=True}). Two points then see the
same bits and the same noise realizations, scaled differently. If $X$ and $Y$ are the estimates
at two neighbouring points,
\[
  \Var(X-Y)=\Var X+\Var Y-2\operatorname{Cov}(X,Y),
\]
and the positive covariance created by shared randomness makes the \emph{differences} between
points much less noisy. The curves look smooth and comparisons between configurations are
fair. Use different seeds only when the points must be statistically independent.
\end{keypoint}

% ------------------------------------------------------------------
\section{Metadata propagation}\label{sec:metadata}
Physical quantities that a downstream block needs but cannot infer from the samples travel in
\code{Signal.meta}. Examples:
\begin{itemize}
  \item the constellation used by the mapper;
  \item the oversampling factor;
  \item the nominal signal power;
  \item the noise variance after equalization;
  \item the accumulated chromatic dispersion of a fiber link.
\end{itemize}
Each block copies the metadata of its input and adds or overrides what it changes.
\Cref{tab:meta-flow} follows the most important keys through a pulse-shaped link. The
mechanism is what lets the \blk{awgn} block set $E_b/N_0$ correctly whatever comes before it,
and lets a soft demapper compute exact LLRs without being told the SNR.
\Cref{app:metadata} lists all keys.

\begin{table}[htbp]
\centering\small
\caption{Metadata along a pulse-shaped QPSK link with $N_{\mathrm{sps}}$ samples per symbol
and $E_s/N_0=\gamma$. The waveform power is per sample; symbols have unit energy.}
\label{tab:meta-flow}
\begin{tabular}{@{}lcccc@{}}
\toprule
after block & \meta{power} & \meta{sps} & \meta{noise_var} & also sets\\
\midrule
\blk{mapper}         & 1 & 1 & -- & \meta{constellation}, \meta{bits_per_symbol}\\
\blk{pulse_shaper}   & $1/N_{\mathrm{sps}}$ & $N_{\mathrm{sps}}$ & -- & \meta{pulse_taps}, \meta{delay}\\
\blk{awgn} (\prm{esn0}) & $1/N_{\mathrm{sps}}$ & $N_{\mathrm{sps}}$ & $N_0=1/\gamma$ & \\
\blk{matched_filter} & 1 & 1 & $1/\gamma$ & \meta{n_symbols}\\
\blk{demapper}       & \multicolumn{4}{l}{uses \meta{constellation} and \meta{noise_var} for exact LLRs}\\
\bottomrule
\end{tabular}
\end{table}

\begin{pitfall}
Metadata describes the \emph{nominal} signal. A block that changes the signal in a way it
cannot describe (a CTLE, an unknown channel) sets the affected keys to \code{None}, and
downstream blocks fall back to measuring. When you write a block, propagate or invalidate
every key your processing changes.
\end{pitfall}

% ------------------------------------------------------------------
\section{Scenarios}\label{sec:scenarios}
A scenario (\file{core/scenario.py}) is a graph plus its run configuration, stored as YAML:
\begin{lstlisting}[style=yaml]
name: QPSK over AWGN
category: Tutorials            # used by the GUI scenario browser
description: ...
seed: 1
batch_size: 64
blocks:
  src:   {type: random_bits, params: {n_bits: 2000}}
  map:   {type: mapper, params: {scheme: psk, order: 4}}
  awgn:  {type: awgn, params: {snr_db: 6}}
  demap: {type: demapper, params: {method: hard}}
  ber:   {type: ber_meter}
  rx:    {type: subsystem, file: subsystems/coded_qam_rx.yaml}   # optional
chain: [src, map, awgn, demap, ber.rx]      # default-port connections
edges:
  - src.out -> ber.tx                        # explicit connections
simulation: {stop: {metric: ber.bit_errors, threshold: 200, max_batches: 500}}
sweep: {params: {awgn.snr_db: [0, 2, 4, 6, 8]}}   # optional: run a sweep instead
probes: [{key: awgn.out, view: constellation}]   # GUI scopes
layout: {src: [0, 0], map: [230, 0]}             # GUI positions
\end{lstlisting}
Files saved from the GUI contain only non-default parameters, so they stay short and readable
and remain usable from the CLI. \Cref{app:scenarios} lists the scenarios shipped with LinkStudio.
